\subsection{Classification closes the $A_8$ fixed sector as the two-qutrit Pauli grading}
Let $V=E_8/(1-g)E_8\cong\mathbb F_3^4$.  The order-three fixed algebra is the
Kac $A_8$ class $\mathfrak a\cong\mathfrak{sl}_9(\mathbb C)$, and the preceding
certificates decompose it into eighty one-dimensional fixed root-orbit lines
\[
 \mathfrak a=\bigoplus_{0\ne v\in V}\mathfrak a_v,
 \qquad [\mathfrak a_v,\mathfrak a_w]\subseteq\mathfrak a_{v+w},
 \qquad \mathfrak a_0=0.
\]
Hence this is a fine $V$-grading.  Bahturin--Kochetov's classification of
Type-$A$ gradings (\texttt{arXiv:0908.0906}) now removes the remaining cocycle
ambiguity up to graded equivalence.  Type II requires a distinguished element
of order two, impossible for $V\cong C_3^4$, so the grading is Type I and comes
from a unique $V$-grading of $M_9(\mathbb C)$.  Restoring the identity gives
all $81=9^2$ matrix components one-dimensional; therefore it is a division
(Pauli) grading, classified by a nondegenerate alternating bicharacter on $V$.
All such forms on $\mathbb F_3^4$ are symplectically equivalent.  Consequently
the fixed $\mathfrak{sl}_9$ grading is, up to symplectic relabeling and
homogeneous rescaling, exactly the tensor two-qutrit Pauli grading studied by
Pelantov\'a--Svobodov\'a--Tremblay
(\texttt{arXiv:quant-ph/0510106}).  Thus the eighty oriented $E_8$ root
triangles are the eighty homogeneous nonidentity Pauli degrees, with W33's
symplectic form deciding exactly which brackets vanish.  The remaining freedom
is normalization/phase choice of one-dimensional generators, not the graded
Lie-algebra class.  No additional physical claim is inferred.
